% AUTO-GENERATED by scripts/build_wp_tables_v11.py -- do not edit by hand.
\begin{table}[htbp]
\centering\small
\caption{\textbf{Attribution (hindsight-formed portfolio).} \emph{Factor attribution, monthly, 36 months.} Factors are built point-in-time from the filing panel: QMJ sorts on ROE, SIZE on total assets (an accounting-size proxy, \emph{not} market capitalisation), MOM on the 12-1 month return, SECTOR is equal-weighted Electric Appliances minus the equal-weighted universe. \textbf{There is no value factor}: a book-to-market sort needs a share count, which this data does not contain point-in-time. Newey--West $t$ (lag 3) in parentheses.}
\label{tab:factors}
\begin{tabular}{lrrrr}
\toprule
& \multicolumn{2}{c}{Portfolio (equal wt.)} & \multicolumn{2}{c}{Mechanical benchmark} \\
\cmidrule(lr){2-3}\cmidrule(lr){4-5}
& \textbf{CAPM${+}$sector} & Five factors & \textbf{CAPM${+}$sector} & Five factors \\
& \scriptsize(headline) & \scriptsize(robustness) & \scriptsize(headline) & \scriptsize(robustness) \\
\midrule
$\alpha$ (annualised) & +12.7\% & +17.8\% & +14.6\% & +23.9\% \\
 & (+2.09) & (+2.29) & (+1.90) & (+3.06) \\
\midrule
MKT & 0.99 & 0.99 & 1.28 & 1.24 \\
 & (+6.76) & (+4.06) & (+6.40) & (+5.11) \\
QMJ & -- & 0.40 & -- & 1.23 \\
 &  & (+0.89) &  & (+2.09) \\
SIZE & -- & 0.66 & -- & 1.09 \\
 &  & (+2.21) &  & (+3.92) \\
MOM & -- & 0.15 & -- & 0.32 \\
 &  & (+0.27) &  & (+1.06) \\
SECTOR & 0.82 & 0.98 & 0.44 & 0.69 \\
 & (+2.82) & (+3.44) & (+1.25) & (+1.79) \\
\midrule
Adjusted $R^2$ & 0.72 & 0.75 & 0.59 & 0.70 \\
\midrule
\multicolumn{5}{l}{\textit{$\alpha$ (annualised) across every specification we ran}} \\
\quad CAPM (1 regressor) & \multicolumn{2}{c}{+12.6\% (+1.54)} & \multicolumn{2}{c}{+14.5\% (+2.04)} \\
\quad CAPM${+}$sector (2 regressors) & \multicolumn{2}{c}{+12.7\% (+2.09)} & \multicolumn{2}{c}{+14.6\% (+1.90)} \\
\quad Four factors (4 regressors) & \multicolumn{2}{c}{+16.0\% (+1.46)} & \multicolumn{2}{c}{+22.6\% (+3.31)} \\
\quad Five factors (5 regressors) & \multicolumn{2}{c}{+17.8\% (+2.29)} & \multicolumn{2}{c}{+23.9\% (+3.06)} \\
\bottomrule
\end{tabular}
\begin{minipage}{\textwidth}\vspace{2pt}\scriptsize
\textit{Notes.} \textbf{The headline specification is CAPM${+}$sector}, chosen on degrees of freedom alone and fixed before the estimates were inspected: 36 monthly observations do not support five regressors. The choice is not self-serving --- the five-factor $\alpha$ is \emph{larger}, so preferring the parsimonious model lowers the $\alpha$ we report. The regressions are run on portfolios formed with 2026 information, so these alphas measure how much in-sample excess return tradable factor exposures span --- not skill. Two features matter for the paper's calibration. The sector loading is large and the portfolio's alpha does \emph{not} vanish once it is included; but the mechanical benchmark, built from the same gate with no discretionary step, earns a \emph{larger} alpha under the headline model (+14.6\% vs.\ +12.7\%) and under every specification in the lower panel. And the quality factor the screen sorts on returned -0.1\% annualised over the sample, consistent with the Layer 1 null. With 36 monthly observations spanning one regime, every alpha here is fragile --- the headline model keeps only two regressors precisely because five is more than this sample supports; see Section~\ref{sec:costs}.
\end{minipage}
\end{table}

